\documentclass[border=10pt]{standalone}
\usepackage[T1]{fontenc}
\usepackage[utf8]{inputenc}
\usepackage{lmodern}
\usepackage[table]{xcolor}
\definecolor{IntelBlue}{HTML}{0068B5}
\definecolor{IntelLightBlue}{HTML}{DDEEF9}
\definecolor{FpgaOrange}{HTML}{E67E22}
\definecolor{FpgaLight}{HTML}{FDEBD0}
\definecolor{HpsGreen}{HTML}{1A7A4A}
\definecolor{HpsLight}{HTML}{D5F5E3}
\definecolor{GrayBg}{HTML}{F2F3F4}
\definecolor{GrayBorder}{HTML}{AAB7B8}
\definecolor{PassGreen}{HTML}{D5F5E3}
\definecolor{PassGreenDark}{HTML}{1E8449}
\definecolor{RowAlt}{HTML}{F7F9FB}
\definecolor{SummaryBg}{HTML}{EBF5FB}
\definecolor{NodeFill}{HTML}{D6EAF8}
\definecolor{NodeBorder}{HTML}{1B4F72}
\definecolor{InitFill}{HTML}{E5E7E9}
\definecolor{InitBorder}{HTML}{717D7E}
\definecolor{SleepFill}{HTML}{FEF5E7}
\definecolor{SleepBorder}{HTML}{935116}
\definecolor{MsgFill}{HTML}{D1F2EB}
\definecolor{MsgBorder}{HTML}{0E6655}
\definecolor{TimeoutRed}{HTML}{C0392B}
\definecolor{BtnBlue}{HTML}{1A5276}
\definecolor{ArrowGray}{HTML}{717D7E}
\usepackage{tikz}
\usetikzlibrary{shapes.geometric, arrows.meta, positioning, fit, backgrounds, calc, decorations.pathreplacing, bending, shadows.blur}

% One legend entry: a short arrow sample drawn in the caller's style, then its
% meaning.
\newcommand{\arrkey}[2]{%
  \tikz[baseline=-0.55ex]{\draw[-{Latex[length=1.7mm,width=1.3mm]},#1]
    (0,0) -- (0.62,0);}~#2}

% Pipeline-stage badge. The SAME numbers are used in three places on purpose:
% the book's Section 4.1 prose, the badges on this figure, and the stage
% comments in fpga_msg_controller.v. A reader can walk 1..5 across all three.
\newcommand{\stg}[1]{%
  \tikz[baseline=-0.6ex]{\node[circle, fill=IntelBlue, text=white,
    font=\tiny\bfseries, inner sep=0.6pt, minimum size=3.7mm] {#1};}}

\begin{document}
% ROUND 2 datapath, synced from the book's Figure 4.1 (project_book_REVISED).
% Adds, versus the round-1 drawing: msg_nav_rom (the category jump table the
% FSM navigates by), msg_text_rom + msg_text_export (message text now lives
% in fabric and crosses the bridge), and the second idle_timer instance (the
% system-idle "sleep" timer, enabled only while msg_nav_rom reports
% in_default). The single "Countdown Timer" of round 1 is relabelled
% "Per-Message Timer" to disambiguate it from the sleep timer.
%
% RESYNC (readability pass): the layout was decluttered to match the book's
% updated Figure 4.1 -- regions I/II/III are now tiled left to right with no
% overlapping dashed boxes, the two full-width "any press" arcs are replaced by
% one tap that feeds both timers on short branches, and the HEX feeds are on
% clean non-crossing lanes.
\begin{tikzpicture}[
  every node/.style={font=\small\sffamily},
  % Block text is deliberately set at \large: in the book the whole picture is
  % shrunk to a fixed 19.5 cm width, so anything smaller becomes unreadable once
  % scaled down. Bigger text relative to the boxes is the single biggest
  % legibility win here.
  pip/.style={rectangle, rounded corners=4pt,
    draw=green!60!black, fill=green!12,
    minimum width=2.4cm, minimum height=1.25cm,
    align=center, font=\large\bfseries},
  romsty/.style={rectangle, rounded corners=4pt,
    draw=MsgBorder, fill=MsgFill,
    minimum width=2.4cm, minimum height=1.25cm,
    align=center, font=\large\bfseries},
  tim/.style={rectangle, rounded corners=4pt,
    draw=red!70!black, fill=red!10,
    minimum width=2.4cm, minimum height=1.25cm,
    align=center, font=\large\bfseries},
  slp/.style={rectangle, rounded corners=4pt,
    draw=SleepBorder, fill=SleepFill,
    minimum width=2.4cm, minimum height=1.25cm,
    align=center, font=\large\bfseries},
  ext/.style={rectangle, rounded corners=4pt,
    draw=gray!60, fill=gray!10,
    minimum width=2.4cm, minimum height=1.25cm,
    align=center, font=\large\bfseries},
  hps/.style={rectangle, rounded corners=4pt,
    draw=blue!70!black, fill=blue!10,
    minimum width=2.4cm, minimum height=2.6cm,
    align=center, font=\large\bfseries},
  bridgebox/.style={rectangle, rounded corners=3pt,
    draw=blue!70!black, fill=white,
    inner xsep=3pt, inner ysep=2pt,
    align=center, font=\footnotesize\itshape\color{blue!70!black}},
  piobank/.style={rectangle, rounded corners=4pt,
    draw=blue!70!black, fill=blue!8,
    minimum width=2.7cm, minimum height=4.6cm,
    align=center, font=\large\bfseries},
  grp/.style={rectangle, rounded corners=6pt,
    draw=black!45, dashed, line width=0.9pt, fill=black!4,
    inner sep=0.42cm},
  % Region headers are the MACRO layer: a reader takes the three numbered
  % regions left to right first, then drops into the blocks inside each. Set
  % bold and large so they read as headers, not footnotes.
  grplbl/.style={font=\large\bfseries\color{IntelBlue},
    inner sep=3pt},
  arr/.style={-{Latex[length=2.5mm,width=1.8mm]},
    draw=black!70, line width=1.0pt},
  % The spine is the path the reader should follow first, so it is the
  % heaviest line on the drawing.
  spine/.style={arr, line width=1.6pt, draw=black!80},
  % Dashed means "this returns to the FSM". Colour still says which subsystem
  % it came from, so the two codings are independent.
  fbk/.style={dash pattern=on 2.4pt off 1.7pt},
  darr/.style={-{Latex[length=2.5mm,width=1.8mm]},
    draw=blue!70!black, line width=1.1pt},
  tarr/.style={-{Latex[length=2.5mm,width=1.8mm]},
    draw=red!80!black, line width=1.2pt},
  navarr/.style={-{Latex[length=2.5mm,width=1.8mm]},
    draw=MsgBorder, line width=1.1pt},
  slparr/.style={-{Latex[length=2.5mm,width=1.8mm]},
    draw=SleepBorder, line width=1.1pt},
  lbl/.style={font=\small\sffamily, fill=white, inner sep=1.5pt},
  % Tap dots mark where a bus is branched rather than broken. At 3.4pt they
  % vanished once the whole picture was scaled down, which made the branches
  % look like they started at the neighbouring block.
  jdot/.style={circle, fill=black!80, inner sep=0pt, minimum size=5.2pt},
  jdotb/.style={jdot, fill=blue!70!black},
  % Step badge. Numbers are taken from the stage comments in
  % fpga_msg_controller.v, so the figure cannot drift from the code.
  badge/.style={circle, fill=IntelBlue, text=white,
    font=\scriptsize\bfseries, inner sep=0.8pt, minimum size=5.4mm},
]

% ============================================================
%  LAYOUT NOTE
%  The picture is spaced generously across the full page width so the
%  \large block text never collides with a badge or an edge label. In the
%  book the whole thing is scaled to a fixed width, so using more width here
%  costs no page height -- it only buys legibility once shrunk. Every arrow
%  runs on its own clear lane; there are no jump-over arcs.
% ============================================================

% ===== Region I: input conditioning (left) =====
\node[ext, minimum width=2.6cm] (keys) at (0.0,  0) {\shortstack{KEY[3:0]\\board pins}};
\node[pip, minimum width=3.0cm] (deb)  at (4.2,  0) {\shortstack{Debouncer\\20\,ms}};
\node[pip, minimum width=3.0cm] (edg)  at (8.0,  0) {\shortstack{Edge\\Detector}};

% ===== Region II: navigation, control and timing (middle) =====
% The FSM is the hub. Its ROMs sit directly beside it (nav above-left, duration
% below-left) and the two timers just outside them, so region II stays a
% compact column that ends well before region III begins.
\node[pip, minimum width=3.8cm, minimum height=1.7cm] (fsm) at (13.4, 0)
  {\shortstack{Message\\FSM}};
\node[romsty, minimum width=3.3cm] (nav) at (12.2,  4.4)
  {\shortstack{Msg Nav ROM\\\normalsize category table}};
\node[slp, minimum width=2.9cm] (slp) at (18.0,  4.4) {\shortstack{Sleep\\Timer}};
\node[romsty, minimum width=3.3cm] (dur) at (12.2, -4.4)
  {\shortstack{Msg Duration\\ROM}};
\node[tim, minimum width=3.3cm] (tmr) at (18.0, -4.4)
  {\shortstack{Per-Message\\Timer}};

% ===== Region III: message text path (right, clearly separated from II) =====
\node[romsty, minimum width=3.0cm] (trom) at (23.0, 0) {\shortstack{Msg Text\\ROM}};
\node[romsty, minimum width=3.3cm] (texp) at (27.6, 0)
  {\shortstack{Text Export\\\normalsize + seqlock}};

% ===== Board display output (sits under the text path, fed by index + seconds)
\node[ext, minimum width=3.0cm] (hex) at (23.0, -4.4)
  {\shortstack{HEX5:HEX0\\\normalsize board 7-seg}};

% ===== PIO register bank + HPS =====
\node[piobank] (pio) at (32.4, 0)
  {\shortstack{PIO\\Register\\Bank\\[3pt]\normalsize read-only}};
\node[hps] (hps) at (36.6, 0) {\shortstack{HPS\\ARM Linux\\LCD app}};

% ---- lookup labels sit beside their lanes, between the FSM and the nav ROM.
\node[lbl, anchor=east] (navlbla) at (12.2, 2.3) {\shortstack[r]{index\\+ action}};
\node[lbl, anchor=west] (navlblb) at (13.7, 2.3) {\shortstack[l]{next\\index}};

% ---- step badges, set OUTSIDE the top-left corner so they never sit on top
%      of the (now larger) block text. Numbers match the stage comments in
%      fpga_msg_controller.v. ROMs carry no badge: they are lookup tables
%      serving a stage, not stages of their own.
\node[badge, anchor=south east] at (deb.north west) {1};
\node[badge, anchor=south east] at (edg.north west) {2};
\node[badge, anchor=south east] at (fsm.north west) {3};
\node[badge, anchor=south east] at (tmr.north west) {4a};
\node[badge, anchor=south east] at (slp.north west) {4b};
\node[badge, anchor=south east] at (hex.north west) {5};

% ---- region grouping: read the three numbered boxes left to right (macro),
%      then look inside each for the blocks (micro).
\begin{scope}[on background layer]
  \node[grp, fit=(keys)(deb)(edg),
        label={[grplbl]above:\textbf{I} \ Input conditioning}] {};
  % The two (10.9,+/-5.6) points pull the region-II box tall enough to enclose
  % the "any press" reload lanes, so the reload routing lives inside region II
  % instead of arcing outside every box the way the old drawing did.
  \node[grp, fit={(nav)(slp)(fsm)(dur)(tmr)(navlbla)(navlblb)(10.9,5.6)(10.9,-5.6)},
        label={[grplbl]above:\textbf{II} \ Navigation, control and timing}] {};
  \node[grp, fit=(trom)(texp),
        label={[grplbl]above:\textbf{III} \ Message text path}] {};
\end{scope}

% ---- fabric / HPS boundary (between the PIO bank and the HPS) ----
\draw[dashed, draw=blue!70!black, line width=1.0pt] (34.5, -7.4) -- (34.5, 5.6);
\node[bridgebox] at (34.5, 4.4) {\shortstack{Lightweight\\HPS--FPGA\\Bridge}};

% ---- main left-to-right spine: one unbroken line, left edge to right edge ----
\draw[spine] (keys.east) -- node[lbl,above]{raw}    (deb.west);
\draw[spine] (deb.east)  -- node[lbl,above]{stable} (edg.west);
\draw[spine] (edg.east)  -- node[lbl,above]{pulse}  (fsm.west);
\draw[spine] (fsm.east)  -- node[lbl,above,pos=0.5]{index}  (trom.west);
\draw[navarr] (trom.east) -- node[lbl,above]{\shortstack{512-bit\\text}} (texp.west);
\draw[darr] (texp.east) -- node[lbl,above]{\shortstack{16 words\\+ status}} (pio.west);
\draw[darr] (pio.east) -- node[lbl,above]{reads} (hps.west);

% ---- FSM <-> nav ROM (category lookup): two clean vertical lanes ----
\draw[navarr]      (12.35, 0.85) -- (12.35, 3.775);
\draw[navarr, fbk] (13.55, 3.775) -- (13.55, 0.85);

% ---- in_default gates the sleep timer ----
\draw[slparr] (nav.east) -- node[lbl,above]{in\_default} (slp.west);

% ---- any button press reloads BOTH timers. One tap on the pulse line just
%      before the FSM feeds two short branches -- up to the sleep timer here,
%      down to the per-message timer below. Both stay inside the region-II box
%      (its fit is stretched to enclose them), so nothing wraps the drawing.
% Tap sits at x=10.2, just LEFT of the ROM left edge (10.55), so the vertical
% branch runs clear of the Msg Nav ROM and Msg Duration ROM boxes.
\node[jdot] at (10.2, 0) {};
\node[lbl, rotate=90] at (10.2, 2.4) {any press};
\draw[slparr] (10.2, 0) -- (10.2, 5.6) -- (18.0, 5.6) -- (slp.north);

% ---- sleep timeout returns to the FSM top face on its own lane ----
\draw[slparr, fbk] (slp.south) -- (18.0, 2.2) -- (14.6, 2.2) -- (14.6, 0.85);
\node[lbl] at (16.4, 2.2) {sleep timeout};

% ---- index down to the duration ROM ----
\draw[arr] (12.2, -0.85) -- (dur.north);

% ---- per-message timer loop: duration ROM -> timer -> back to the FSM ----
\draw[tarr] (dur.east) -- node[lbl,above]{duration} (tmr.west);
\draw[tarr, fbk] (tmr.north) -- (18.0, -2.2) -- (14.6, -2.2) -- (14.6, -0.85);
\node[lbl] at (16.4, -2.2) {msg timeout};
% down branch of the same "any press" tap: reload the per-message timer, then
% carry the identical edge-detector pulse on to HEX2 (last button pressed),
% so one signal serves both instead of a second full-width arc.
\draw[tarr] (10.2, 0) -- (10.2, -5.6) -- (17.5, -5.6) -- (17.5, -5.025);
\node[jdot] at (17.5, -5.6) {};
\draw[tarr] (17.5, -5.6) -- (22.0, -5.6) -- (22.0, -5.025);
\node[lbl, below] at (19.9, -5.6) {last button};

% ---- index also drives the HEX digits (message number on the 7-seg), tapped
%      from the spine in the clear gap between regions II and III ----
\node[jdot] at (20.3, 0) {};
\draw[arr] (20.3, 0) -- (20.3, -1.6) -- node[lbl,above,pos=0.8]{index} (23.0, -1.6) -- (hex.north);

% ---- state + index up to the PIO bank, on a clear lane that rides above the
%      region-III header so the two never touch ----
\draw[darr] (15.2, 0.85) -- (15.2, 2.05) -- ($(pio.west)+(0,2.05)$);
\node[lbl, above] at (28.5, 2.05) {state + index};

% ---- seconds remaining under the fabric to the PIO bank ----
\draw[darr] ([xshift=6mm]tmr.south) -- (18.6, -6.8) -- (32.4, -6.8) -- (pio.south);
\node[lbl, below] at (27.4, -6.8) {seconds + timeout};
\node[jdotb] at (23.0, -6.8) {};
\draw[darr] (23.0, -6.8) -- node[lbl,right]{seconds} (hex.south);

% ---- legend, drawn as a node in the SAME coordinate system as the picture
%      above (rather than as ordinary paragraph text after \end{tikzpicture})
%      so it is reliably centered under the diagram regardless of the
%      surrounding document's text width.
\node[align=center, font=\footnotesize, text width=27cm] at (16.5, -8.6) {%
\arrkey{draw=black!80, line width=1.6pt}{forward datapath}\hspace{0.55cm}%
\arrkey{draw=MsgBorder, line width=1.1pt}{navigation and message text}\hspace{0.55cm}%
\arrkey{draw=red!80!black, line width=1.2pt}{per-message timer}\\[3pt]
\arrkey{draw=SleepBorder, line width=1.1pt}{sleep path}\hspace{0.55cm}%
\arrkey{draw=blue!70!black, line width=1.1pt}{read by the HPS}\hspace{0.55cm}%
\arrkey{draw=black!70, line width=1.1pt, dash pattern=on 2.4pt off 1.7pt}{dashed: returns to the FSM}\\[3pt]
\stg{1}\,--\,\stg{5}~pipeline stage, numbered as in the book text and in
\texttt{\small fpga\_msg\_controller.v}%
};

\end{tikzpicture}

\end{document}
